\section{Methodology}

\subsection{Research Design}

The design set out in this section tests the five hypotheses of Section 2. This study follows the empirical strategy of \textcite[Section 4]{beirne_sugandi_2023} and modifies it in documented ways. The object of estimation is a country-specific structural VAR for Japan on working-days data, organized in two tiers. Tier one is the daily restricted VAR over the paper period, the specification that produces the paper's published headline figures. Tier two is the monthly VAR over the paper period and the extended sample, which matches the paper's appendix structure and provides the frequency robustness check. The paper period runs from 14 January 1999 to 31 March 2021, and the extended sample runs through 30 June 2026.

The scope of the replication is the Japanese model. The paper estimates its VARs for three safe haven destinations, and this study re-estimates the Japanese model only, with the Swiss and United States estimates entering the discussion only as the published evidence reports them. The paper's weekly frequency panels are likewise not replicated, and the two-tier design covers the daily and monthly frequencies only. Both exclusions are scope decisions, disclosed so that the boundary of the replication claim is explicit.

The working-days calendar is constructed by removing weekends from the full daily calendar, keeping only Monday through Friday. The sample start is the paper's first working day on which all financial series are available. The daily estimation sample contains 5,199 complete-case observations in the paper period and 1,234 in the extension, for a full-sample total of 6,433 working days. The monthly tier contains 264 complete months for the paper period and 63 for the extension. The risk-off variable is anchored to the VIX, which trades on every United States market day, and a left join preserves dates on which the VIX exists but the Japanese series do not, such as Japanese holidays, so that risk-off episode start dates falling on Japanese non-trading days are retained.

\subsection{Data and Sources}

Two groups of sources supply the data for the design. The daily financial market series were downloaded from LSEG Workspace. These are the VIX, the Nikkei 225 index, the bilateral USD/JPY exchange rate, the 10-year JGB yield, and the 10-year US Treasury yield. All LSEG exports share a six-column format with the date and open, high, low, close, and volume fields, so a single parser handles them. The JGB file is the exception in layout, with a title row that is skipped and the 10-year yield stored in the eleventh column, and rows without a valid 10-year value are dropped. The spread variable is computed as the 10-year JGB yield minus the 10-year US Treasury yield.

The macroeconomic series were collected via automated Python API scripts from their official sources. Real GDP and nominal GDP come from the Federal Reserve Bank of St.\ Louis, FRED series JPNRGDPEXP and JPNNGDP, retrieved in 2026. The World Uncertainty Index (WUI) for Japan comes from the World Uncertainty Index database maintained by Ahir, Bloom, and Furceri, retrieved through FRED as series WUIJPN at the quarterly frequency. The real effective exchange rate comes from the BIS broad index. The capital flow series come from the Bank of Japan's BPM6 balance of payments statistics, measuring the net incurrence of portfolio investment liabilities, other investment liabilities, and direct investment liabilities.

The capital flow extraction is the most demanding part of the data preparation. Pre-2014 observations come from the Bank of Japan's 6pi-1 portfolio summary, a Shift-JIS encoded CSV with a mixed layout that requires reading the raw lines, locating the monthly figures section, and parsing a 29-column table in which the equity, long-term debt, and short-term debt net liability columns sit at fixed positions. The values are comma-formatted and use a dash placeholder for missing observations, and the dates are written in Japanese era format, which is converted to the Gregorian calendar with the Heisei era offset of 1988 and the Reiwa offset of 2018. Post-2014 observations come from the Bank of Japan's BPM6 API files, which are simple two-column CSVs in units of 100 million yen. The debt securities series is the sum of the long-term and short-term debt securities components, consistent with the Bank of Japan's D-04 reporting convention, and the pre-2014 and post-2014 segments are concatenated into continuous series for each flow category.

Table~\ref{tab:variables} defines each variable, its frequency, and its source. The vintage arrangements for the paper period are described in the next subsection.

\begin{table}[H]
  \centering
  \caption{Endogenous Variables, Definitions, and Sources}
  \label{tab:variables}
  \small
  \begin{tabularx}{\textwidth}{@{}p{2.3cm} X p{2.5cm} p{3.7cm}@{}}
    \toprule
    Variable & Definition & Frequency & Source \\
    \midrule
    Risk-off & Indicator equal to 1 when the VIX exceeds its 60-day moving average by at least 10 points & Daily & LSEG Workspace, authors' calculation \\
    Log (WUI) & Log of the World Uncertainty Index for Japan & Quarterly to daily & World Uncertainty Index database \\
    Spread & 10-year JGB yield minus 10-year US Treasury yield & Daily & LSEG Workspace \\
    Log (RGDP) & Log of constant-price real GDP & Quarterly to daily & FRED ALFRED (vintage 2021-06-01) \\
    Log (REER) & Log of the real effective exchange rate, broad index & Monthly to daily & BIS (Wayback 2021-05-24) \\
    Log (Nikkei) & Log of the Nikkei 225 index & Daily & LSEG Workspace \\
    Debtsec & Net portfolio investment inflows to debt securities, percent of nominal GDP & Monthly to daily & Bank of Japan \\
    Equity & Net portfolio investment inflows to equity, percent of nominal GDP & Monthly to daily & Bank of Japan \\
    Other & Net other investment inflows, percent of nominal GDP & Monthly to daily & Bank of Japan \\
    Direct & Net direct investment inflows, percent of nominal GDP & Monthly to daily & Bank of Japan \\
    \bottomrule
  \end{tabularx}
  \caption*{\textit{Note.} Capital flow variables follow BPM6 conventions. Positive values indicate capital inflows. The nominal GDP denominator is current-vintage throughout.}
\end{table}

\subsection{Dual-Vintage Data Policy}

The data follow a dual-vintage policy because the published sample has been retroactively revised since 2021. Japan's real GDP underwent a base-year change from chained 2011 yen to chained 2015 yen in December 2020, and subsequent benchmark revisions raised current-vintage output by 3.5 to 6.3 percent over 2019 to 2021 while flattening the COVID-19 trough. The BIS rebased its REER from 2010 equal to 100 to 2020 equal to 100 in January 2023, revising the entire historical series and trade weights.

The paper period therefore uses the real GDP vintage of 1 June 2021 from the FRED ALFRED archive, retrieved through the FRED API with the vintage date parameter set to 2021-06-01. The REER for the paper period uses the BIS broad index snapshot of 24 May 2021, recovered through the Internet Archive Wayback Machine from the archived BIS file, because the BIS does not publicly archive retrospective vintages and the Wayback snapshot is the closest approximation of the file the original authors observed. The extended period uses current-vintage data for all series. Replication requires the data as the original authors observed them, while extension requires the best available current estimates, and the two goals impose incompatible data demands that the dual-vintage policy resolves explicitly.

\subsection{Variables}

The model contains ten endogenous variables in the Cholesky ordering of the paper, all drawn from the series described above. The risk-off indicator enters first because it is the exogenous shock of interest. It is followed by the log of the WUI, the JGB yield spread relative to the 10-year US Treasury yield, the log of real GDP, the log of the REER, the log of the Nikkei 225 index, and the four capital flow variables expressed as percentages of nominal GDP, as defined in Table~\ref{tab:variables}.

A risk-off event is defined following \textcite{debock_carvalhofilho_2013} as a day on which the VIX exceeds its 60-day backward-looking moving average by 10 percentage points, and the daily risk-off variable is a binary indicator for such days. This definition reproduces all fifteen in-sample episode initial dates from the paper's Table 1 exactly to the day. At the monthly frequency the daily binary is aggregated as the proportion of risk-off days in the month, which preserves the signs of the impulse responses, whereas summing the daily values inflates the shock standard deviation to 2.23 and produces excessively sharp response curves, and a binary indicator for any risk-off day in the month reverses the output sign. The paper does not state its monthly aggregation method.

All level variables are transformed before estimation. The stock index, the REER, and real GDP enter as natural logs. The WUI enters as a natural log with a positivity guard, since early observations can reach zero and the log of zero is undefined, so non-positive observations are set to missing before the transform. The capital flow variables are converted from their native units into percentages of nominal GDP in three steps. The monthly flows in 100 million yen are converted to billions of United States dollars using the end-of-month USD/JPY rate, the last trading day observation of each month, with the conversion multiplying by 0.1 and dividing by the rate. Nominal GDP in millions of yen is converted to billions of United States dollars with the same end-of-month rate, dividing by the rate times 1,000, and the quarterly GDP is expanded to a full monthly sequence with the quarterly value carried forward so that the second and third months of each quarter are not missing. The percentage of GDP is then the flow in dollars divided by GDP in dollars, multiplied by 100. Positive values therefore indicate capital inflows under the BPM6 net incurrence convention.

The exogenous variables are a constant, a linear time trend, and eleven monthly seasonal indicators with January as the reference month. The paper states that seasonal and time dummies enter the model without specifying their form. Monthly indicators are the natural reading for a daily-frequency VAR in which the lower-frequency series enter through interpolation, because the interpolation injects artificial intra-period variation that monthly indicators absorb. The linear time trend is indexed from 1 to the number of observations.


\subsection{Econometric Model}

These variables enter a VAR whose generic form follows \textcite{beirne_sugandi_2023}, with the lag length selected by information criteria, as stated in Equation~\eqref{eq:var},

\begin{equation}
  Y_t = \sum_{\tau = 1}^{k} Y_{t - \tau} A_\tau + X_t B + c + \varepsilon_t,
  \label{eq:var}
\end{equation}

where \(Y_t\) is the vector of the ten endogenous variables defined in Equation~\eqref{eq:vector}, \(X_t\) is the matrix of exogenous variables, \(A_\tau\) and \(B\) are coefficient matrices, \(c\) is the vector of constants, \(\varepsilon_t\) is the vector of reduced-form residuals, and \(k\) is the lag length. The multiplication order follows the paper's notation, in which the lagged vectors are post-multiplied by the coefficient matrices, and the scalar form of each equation is written out below, in Equations~\eqref{eq:riskoff} and~\eqref{eq:rgdp}.

In the estimated model the endogenous vector, Equation~\eqref{eq:vector}, contains the ten actual variables in their Cholesky order,

\begin{equation}
  Y_t = \begin{pmatrix} \text{risk\_off}_t & \text{log\_wui}_t & \text{spread}_t & \text{log\_rgdp}_t & \text{log\_reer}_t \\
  \text{log\_nikkei}_t & \text{debtsec\_pct}_t & \text{equity\_pct}_t & \text{other\_pct}_t & \text{direct\_pct}_t \end{pmatrix}',
  \label{eq:vector}
\end{equation}

where risk\_off is the binary risk-off indicator, log\_wui is the log of the World Uncertainty Index, spread is the 10-year JGB yield minus the 10-year US Treasury yield, log\_rgdp is the log of constant-price real GDP, log\_reer is the log of the real effective exchange rate, log\_nikkei is the log of the Nikkei 225 index, and the four flow variables are the net incurrence of portfolio debt, portfolio equity, other, and direct investment liabilities as percentages of nominal GDP. The exogenous matrix contains the constant, the linear time trend, and the eleven monthly indicators, so the \(j\)-th row is \(x_t' = (1, t, d_{2t}, \ldots, d_{12t})\). The lag length \(k\) is 20 in the daily tier and 2 or 3 in the monthly tier, as documented below.

The block exogeneity restriction is imposed on the first equation of the system. The risk-off indicator is allowed to depend only on its own lagged values and the exogenous variables,

\begin{equation}
  \text{risk\_off}_t = c_1 + \sum_{\tau = 1}^{k} a_{1,1,\tau}\, \text{risk\_off}_{t-\tau} + x_t' b_1 + \varepsilon_{1,t},
  \label{eq:riskoff}
\end{equation}

with the coefficients on the lagged values of the other nine endogenous variables in this equation restricted to zero. Every other equation depends on the lagged values of all ten variables. The real GDP equation, which produces the headline result, is written out in Equation~\eqref{eq:rgdp},

\begin{equation}
\begin{split}
  \text{log\_rgdp}_t = c_4 + & \sum_{\tau = 1}^{k} a_{4,1,\tau}\, \text{risk\_off}_{t-\tau} + \sum_{\tau = 1}^{k} a_{4,2,\tau}\, \text{log\_wui}_{t-\tau} + \cdots \\
  & + \sum_{\tau = 1}^{k} a_{4,10,\tau}\, \text{direct\_pct}_{t-\tau} + x_t' b_4 + \varepsilon_{4,t},
\end{split}
  \label{eq:rgdp}
\end{equation}

and the remaining eight equations share the same structure with the respective dependent variable. The unrestricted specification removes the restriction in Equation~\eqref{eq:riskoff} and lets the risk-off equation depend on the lagged values of all ten variables, and it is reported as a robustness check mirroring the paper's appendix structure.

Identification is recursive. The reduced-form residuals are orthogonalized with a lower-triangular Cholesky decomposition of the residual covariance matrix, and the impulse responses are generated from a one-standard-deviation structural shock to the risk-off variable, which is ordered first because it is the exogenous shock of interest. Confidence intervals are reported at the 95 percent level.

The restricted specification is estimated by ordinary least squares equation by equation with a custom implementation, and the unrestricted specification uses the statsmodels VAR implementation. The impulse response horizon is 125 trading days for the daily tier, approximately six months, and 40 months for the monthly tier. All computations run in Python with the statsmodels library, and the complete pipeline is the single script \texttt{scripts/replicate\_svar.py} in the public repository for this study. The random number generator is seeded at 42 for the Monte Carlo procedure.

The lag length is selected asymmetrically across tiers, after four candidate specifications were tested and eliminated in sequence. The first attempt used the Akaike Information Criterion (AIC) for both tiers, following the paper's stated method. The daily AIC declined monotonically through every lag from 1 through 20 and selected the boundary, and extending the search to 30 lags produced the same behavior, because each additional lag adds 100 endogenous coefficients in a ten-variable system with over 5,000 observations and the likelihood improvement from absorbing residual autocorrelation swamps the AIC penalty of 2 per parameter. The second attempt applied the AIC to the monthly tier, and the same boundary selection appeared through the maximum of 12. The third attempt switched the monthly tier to the Bayesian Information Criterion (BIC), whose penalty of roughly 5.7 per parameter at the monthly frequency produced interior minima instead of boundary selection, at two lags for the paper period and three lags for the extended sample, and the two-lag paper-period specification reproduces the paper's monthly figures. Lag 4 widens the confidence intervals and adds a second oscillation absent from the paper's plots. The fourth attempt tried the BIC on the daily tier, which selected one or two lags and produced responses too persistent to reproduce the paper's daily decay rates, so it was rejected. The resolution is asymmetric. The daily tier retains the AIC with a maximum lag of 20, following the paper's stated method with no better alternative, and the monthly tier uses the BIC with a maximum lag of 12, returning two lags for the paper period and three lags for the extended sample.

The confidence intervals adopt different methods by tier. The restricted model's intervals use a Monte Carlo delta method with 1,000 draws. Each draw samples synthetic residuals from a normal distribution with the estimated residual covariance matrix, reconstructs a bootstrap data path from the estimated coefficient matrices, re-estimates the VAR on the synthetic path, recomputes the impulse responses, and orthogonalizes them with the original Cholesky factor, and the 2.5th and 97.5th percentiles of the 1,000 replications form the band. The unrestricted model's intervals use asymptotic standard errors from the statsmodels impulse response analysis with orthogonalized errors, the same method EViews applies by default. Both interval types widen with the forecast horizon, which distinguishes them from the paper's published bands, which appear flat.

The model is estimated in levels, which is standard in the macro VAR literature \parencite{sims_1980,lutkepohl_2005}. Augmented Dickey-Fuller tests reported in Section 4 confirm that the risk-off indicator, the uncertainty index, and the four capital flow variables are stationary, and that the spread, real GDP, the REER, and the stock index do not reject a unit root at conventional significance levels, consistent with trend-stationary or slowly moving levels behavior that the levels specification does not depend on.

\subsection{Data Transformation Pipeline}

The estimation dataset for the specifications above passes through three transformation stages. In the first stage, each series is stored at its native frequency, and the raw files are re-read without the sample start filter so that the full available history precedes 14 January 1999. These pre-sample observations serve as anchor points for the interpolation, because a quadratic fitted to the in-sample observations alone would have nothing to anchor its boundary behavior.

In the second stage, the lower-frequency series are interpolated to the daily trading-day grid using quadratic-match average interpolation. The method fits a local quadratic through each triplet of adjacent source observations, constrained so that the mean of the interpolated high-frequency values within each low-frequency period equals the source observation. This reimplements the documented EViews default for frequency conversion, which the paper does not name, and the mean-preservation property is verified at machine precision across all series. The interpolation runs over the trading-day grid directly rather than over calendar days with subsequent subsetting, because the quadratic polynomial differs across grids, and it does not extrapolate beyond the last source observation, so trailing missing values for the quarterly series are forward-filled to carry them through the end of the sample. The interpolation is implemented in Python with pandas, because R has no native quadratic spline with the required boundary behavior.

In the third stage, the daily and monthly estimation datasets are assembled. The daily dataset joins all ten variables on the date in Cholesky order and filters to the working-days calendar, with the risk-off indicator, the spread, and the log stock index entering directly from the daily financial series and the remaining seven variables entering from the interpolated series. The monthly dataset takes the last trading day observation of each month for the financial variables and the proportion of risk-off days for the risk-off indicator. The real GDP and REER series for the paper period are overwritten with the vintage data, with the ALFRED vintage GDP assigned to each month within its quarter and the vintage REER assigned to its calendar month. The exogenous matrix stacks the constant, the linear time trend, and the eleven monthly indicators, and the same matrix enters both the daily and monthly specifications.

\subsection{Contrast with the Paper's Methods and Unreported Settings}

The implementation choices made above are contrasted with the paper's methods in Table~\ref{tab:contrast}. Every row in which the paper is silent on a setting records what had to be reconstructed and the basis for the reconstruction. The paper states its frequency structure, its Cholesky ordering, and its recursive restriction. It does not state the quadratic interpolation variant, the maximum lag searched, the form of the seasonal and time dummies, the confidence interval method, or the monthly aggregation of the risk-off indicator. Each of these settings was reverse-engineered by testing the candidates against the paper's published figures and keeping the version that reproduced them, with the divergence disclosed in the table.

\begin{table}[htbp]
  \centering
  \caption{Contrast with the Paper's Methods and Unreported Settings}
  \label{tab:contrast}
  \small
  \begin{tabularx}{\textwidth}{@{}p{2.2cm} p{3.6cm} p{5.0cm} X@{}}
    \toprule
    Aspect & Paper & This replication & Basis \\
    \midrule
    Frequency & Daily baseline, monthly appendix & Daily tier 1, monthly tier 2 & Paper's own structure \\
    Model & Restricted VAR (block exogeneity) & Restricted tier 1, unrestricted tier 2 & Paper reports both forms, the replication compares them \\
    Lag selection & AIC, maximum unreported & AIC daily (maxlag 20), BIC monthly (lag 2 paper, lag 3 extended) & AIC overfits to the boundary lag at monthly frequency \\
    Identification & Cholesky, risk-off first & Same ordering & Paper stated \\
    Interpolation & EViews quadratic, variant unreported & Python quadratic-match average & Documented EViews default, mean preservation verified \\
    Seasonal and time dummies & Stated, form unreported & Eleven monthly indicators plus linear trend & Natural reading for daily data with interpolated series \\
    Confidence intervals & Method unreported & Monte Carlo delta (restricted), asymptotic (unrestricted) & Both widen with the horizon, unlike the paper's published bands \\
    Monthly risk-off aggregation & Unreported & Proportion of risk-off days & Summing inflates the shock variance, a binary flips the output sign \\
    Capital flows & IMF IFS quarterly & Bank of Japan BPM6 monthly & BPM6-comparable, higher frequency \\
    Data vintage & Mid-2021 downloads & ALFRED 2021-06-01 GDP and Wayback 2021-05-24 REER for the paper period & Benchmark revisions retroactively change historical values \\
    Software & EViews, inferred from figure styling & Python, statsmodels, custom OLS & Independent reimplementation \\
    Sample & 14 January 1999 to 31 March 2021 & Same, extended through 30 June 2026 & Generalizability test \\
    \bottomrule
  \end{tabularx}
\end{table}

The reverse engineering was not arbitrary. For the interpolation, the paper names quadratic interpolation without a variant, and quadratic-match average is the documented EViews default for converting between frequencies, so it was adopted and its mean-preservation property verified at machine precision. For the lag selection, the AIC overfits to the boundary lag at both frequencies on current data, and the monthly tier was switched to the BIC, which selects two lags for the paper period and three for the extended sample and reproduces the paper's monthly figures. For the risk-off aggregation, summing the daily values inflated the shock standard deviation to 2.23 and produced excessively sharp impulse responses, a binary indicator for any risk-off day in the month reversed the output sign, and the proportion of risk-off days preserved the signs and matched the paper's appendix figures. Each choice is therefore documented as a test outcome rather than a preference.

\subsection{Replication Integrity}

The replication initially produced impulse responses that did not match the paper's published figures, and the failure turned out to be diagnostic rather than terminal. Five code defects were identified and corrected. The first was a transpose error in the impulse response recursion that swapped equation and variable indices, corrupting every impulse response in every specification. The second was a plot indexing error that displayed the response of the risk-off variable to shocks in other variables instead of the response of each variable to a risk-off shock. The third was a missing constant term, which biased the coefficients in a system where several variables have non-zero means. The fourth was flat confidence bands built from residual standard errors that did not propagate through the impulse response recursion. The fifth was a coefficient storage bug confined to the restricted specification, which misaligned the own-lag coefficients of the risk-off equation and inflated the apparent similarity between the paper period and the extended period.

A second finding emerged after the code defects were corrected. On current-vintage data the restricted model's real GDP response to a risk-off shock had the wrong sign, and substituting the ALFRED 2021-06-01 vintage restored the paper's negative, persistent, and statistically significant response. The sign reversal is attributable to Japan's post-2022 GDP benchmark revisions, which raised current-vintage output over 2019 to 2021 and altered the COVID-19 trough. This finding is the empirical justification for the dual-vintage policy, and it is examined further in the discussion.

\subsection{Limitations of the Methodology}

The design carries limitations that bound the interpretation of the results. The paper does not report which quadratic-match variant EViews used, what lag maximum was searched, or how its confidence intervals were computed, so every independent replication of the paper carries irreducible uncertainty about whether differences originate in data, code, or configuration. The BIS does not publicly archive retrospective REER vintages, so the vintage REER rests on a single Wayback Machine snapshot. The WUI is available at quarterly frequency for the full sample, while the paper states a monthly series that cannot exist before 2008. The capital flow variables come from the Bank of Japan's BPM6 statistics rather than the paper's IMF IFS source, and the two are directly comparable under BPM6 conventions but differ in frequency and in the details of component classification. The Monte Carlo confidence bands differ between processor architectures at machine precision, with differences in the sixth decimal place that do not affect the match classifications but would prevent bit-for-bit reproduction across hardware. These limitations are addressed in the discussion alongside their consequences for specific results. Section 4 reports the results produced by this design, from the descriptive statistics and stylized facts through the impulse responses and the cross-validation against the published paper.
